\documentclass[12pt, a4paper]{article}

% ---- Packages ----
\usepackage[T1]{fontenc}
\usepackage[utf8]{inputenc}
\usepackage{amsmath, amssymb}
\usepackage{booktabs}
\usepackage{geometry}
\usepackage{hyperref}
\usepackage{xcolor}
\usepackage{array}
\usepackage{caption}
\usepackage{longtable}

\geometry{margin=2.5cm}

% ---- Title ----
\title{\textbf{二元体系相平衡浓度推导}\\[0.5em]
\large 热力学 / 材料科学笔记}
\author{}
\date{}

\begin{document}

\maketitle
\tableofcontents
\newpage

%% ==============================================================
\section{已知条件与符号定义}
%% ==============================================================

考虑一个包含两个组分（记为 elem~1 和 elem~2）的体系，其物理参数定义如下。

\subsection*{初始状态}
\begin{itemize}
  \item 初始浓度：$C_1^0$（组分~1），$C_2^0$（组分~2）
  \item 初始温度：$T_0$
\end{itemize}

\subsection*{物理参数}
\begin{itemize}
  \item 液相线斜率：$m_1$（组分~1），$m_2$（组分~2）
  \item 固液分配系数：$k_1$（组分~1），$k_2$（组分~2）
  \item 温度修正项：$\Delta T_R$（已知常数）
\end{itemize}

\subsection*{当前状态（在某一特定温度 $T$ 时）}
\begin{itemize}
  \item 液相实际浓度：$C_1$，$C_2$
\end{itemize}

\subsection*{求解目标}
\begin{itemize}
  \item 平衡浓度：$C_1^{eq}$，$C_2^{eq}$
  \item 固相分数变化量：$\Delta f_S$
\end{itemize}

%% ==============================================================
\section{基础方程构建}
%% ==============================================================

根据热力学线性近似及溶质再分配原理，建立以下两个核心方程：

\paragraph{方程~(1)：温度–浓度线性关系}

假设温度变化由各组分浓度变化引起的液相线移动叠加而成，并包含修正项 $\Delta T_R$：
\begin{equation}
  T = T_0 + m_1\!\left(C_1^{eq} - C_1^0\right) + m_2\!\left(C_2^{eq} - C_2^0\right) + \Delta T_R
  \label{eq:temp}
\end{equation}

\paragraph{方程~(2)：固相分数与分配系数关系}

基于杠杆定律及分配系数定义，固相分数变化 $\Delta f_S$ 在两个组分中应保持一致：
\begin{equation}
  \Delta f_S
  = \frac{C_1^{eq} - C_1}{C_1^{eq}(k_1 - 1)}
  = \frac{C_2^{eq} - C_2}{C_2^{eq}(k_2 - 1)}
  \label{eq:fs}
\end{equation}

%% ==============================================================
\section{推导过程}
%% ==============================================================

\subsection{步骤~1：利用方程~(2) 建立 $C_1^{eq}$ 与 $C_2^{eq}$ 的关系}

由方程~\eqref{eq:fs} 可得：
\[
  \frac{C_1^{eq} - C_1}{C_1^{eq}(k_1 - 1)}
  = \frac{C_2^{eq} - C_2}{C_2^{eq}(k_2 - 1)}
\]

整理后：
\[
  C_2^{eq}(k_2 - 1)(C_1^{eq} - C_1)
  = C_1^{eq}(k_1 - 1)(C_2^{eq} - C_2)
\]

令 $A_1 = k_1 - 1$，$A_2 = k_2 - 1$，则：
\[
  A_2\,C_2^{eq}\,C_1^{eq} - A_2\,C_2^{eq}\,C_1
  = A_1\,C_1^{eq}\,C_2^{eq} - A_1\,C_1^{eq}\,C_2
\]

解出 $C_2^{eq}$ 关于 $C_1^{eq}$ 的表达式：
\begin{equation}
  C_2^{eq}
  = \frac{A_1\,C_2\,C_1^{eq}}{A_2\,C_1^{eq} - (A_2 - A_1)\,C_1^{eq} + A_1\,C_1^{eq} - A_2\,C_1}
  \quad \Longrightarrow \quad
  C_2^{eq} = \frac{A_1\,C_2\,C_1^{eq}}{A_1\,C_1^{eq} - A_2(C_1^{eq} - C_1) + A_2\,C_1^{eq} - A_1\,C_1^{eq}}
  \label{eq:C2eq_inter}
\end{equation}

更简洁地写为：
\begin{equation}
  C_2^{eq} = \frac{(k_1-1)\,C_2\,C_1^{eq}}{(k_2-1)\,C_1^{eq} - \bigl[(k_2-1)C_1 - (k_1-1)C_1\bigr]}
  \label{eq:C2eq}
\end{equation}

\subsection{步骤~2：代入方程~(1) 求解}

将 $C_2^{eq}$ 代入温度方程~\eqref{eq:temp}，令 $\Delta T = T - T_0 + m_1 C_1^0 + m_2 C_2^0 - \Delta T_R$。

分母部分化简为：
\[
  D = (k_2 - 1)\,C_1^{eq} + (k_1 - k_2)\,C_1
\]

则：
\[
  C_2^{eq} = \frac{(k_1-1)\,C_2\,C_1^{eq}}{D}
\]

代回主方程并移项，整理得到关于 $C_1^{eq}$ 的一元二次方程：
\begin{equation}
  a\,\left(C_1^{eq}\right)^2 + b\,C_1^{eq} + c = 0
  \label{eq:quadratic}
\end{equation}

%% ==============================================================
\section{最终结果}
%% ==============================================================

\subsection{二次方程系数}

定义关于 $C_1^{eq}$ 的系数如下：

\begin{equation}
  a = m_1(k_2-1) + m_2(k_1-1)\frac{C_2}{C_1}
  \label{eq:coef_a}
\end{equation}

\begin{equation}
  b = m_1\bigl[(k_1-k_2)C_1\bigr] + m_2(k_1-1)C_2 - \Delta T\,(k_2-1)
  \label{eq:coef_b}
\end{equation}

\begin{equation}
  c = -\Delta T\,(k_1-k_2)\,C_1
  \label{eq:coef_c}
\end{equation}

其中：
\[
  \Delta T = T - T_0 - m_1 C_1^0 - m_2 C_2^0 - \Delta T_R
\]

\subsection{求解 $C_1^{eq}$}

\begin{equation}
  C_1^{eq} = \frac{-b + \sqrt{b^2 - 4ac}}{2a}
  \label{eq:C1eq}
\end{equation}

\textbf{注：} 根据物理意义，浓度必须为正（凝固过程中通常大于初始浓度），故取正根。若 $a \to 0$（即 $k_1 \approx k_2$），则退化为线性方程求解。

\subsection{求解 $C_2^{eq}$}

\begin{equation}
  C_2^{eq} = \frac{(k_1-1)\,C_2\,C_1^{eq}}{(k_2-1)\,C_1^{eq} + (k_1-k_2)\,C_1}
  \label{eq:C2eq_final}
\end{equation}

\subsection{求解 $\Delta f_S$}

\begin{equation}
  \Delta f_S = \frac{C_1^{eq} - C_1}{C_1^{eq}(k_1-1)}
  \label{eq:DeltaFS}
\end{equation}

\bigskip
\noindent\textit{注：本文档由 LaTeX 源码转换生成。}

%% ==============================================================
\newpage
\section{CoCrNi-Ti-C 系统参数提取}
%% ==============================================================

\textit{本节参数来源：《3D Cellular Automaton Simulation of Dendrite Growth in CoCrNi-Ti-C Alloys: A Comprehensive Parameterization Framework》}

本节参数用于填充上述二元体系相平衡浓度推导公式中的 $m_1$、$m_2$、$k_1$、$k_2$ 等参数。

\subsection{系统设定}

本文档将 CoCrNi-Ti-C 体系处理为伪二元/伪三元系统：
\begin{itemize}
  \item elem~1 = 钛 (Ti)
  \item elem~2 = 碳 (C)
  \item 基体 = 等原子比 CoCrNi（作为有效溶剂 effective matrix）
\end{itemize}

\subsection{热力学参数}

\begin{table}[h!]
\centering
\caption{热力学参数（$m_1$、$m_2$、$k_1$、$k_2$）}
\begin{tabular}{lllllc}
\toprule
符号 & 参数含义 & Ti (elem 1) & C (elem 2) & 单位 & 数据来源 \\
\midrule
$k_1$ & 平衡溶质分配系数 & $0.06 \sim 0.12$ & $0.03 \sim 0.08$ & 无量纲 & CALPHAD (TCHEA6); [9][15] \\
$m_1$ & 液相线斜率 & $-20 \sim -35$ & $-40 \sim -80$ & K/at.\% & CALPHAD; [9][13][15] \\
$T_0$ & CoCrNi 基体液相线温度 & \multicolumn{2}{c}{1610 K (1337 °C)} & K & 实验 DSC + CALPHAD [10] \\
$T_\text{solidus}$ & 平衡固相线温度 & \multicolumn{2}{c}{1443 K (1170 °C)} & K & 实验 DSC + CALPHAD [10] \\
\bottomrule
\end{tabular}
\end{table}

\noindent\textbf{说明：}
\begin{itemize}
  \item $k < 1$：Ti 和 C 均强烈偏析到液相
  \item $m < 0$：溶质富集导致液相线温度下降
  \item 精确值需通过 Thermo-Calc + TCHEA6 数据库在具体成分的等温截面上提取
\end{itemize}

\subsection{$\Delta T_R$ 修正项}

对于特定成分的 $(\text{CoCrNi})_{1-2x}\text{Ti}_x\text{C}_x$ 合金，初始液相线温度相较纯 CoCrNi 基体发生下移：
\[
  \Delta T_R = m_1 \cdot x + m_2 \cdot x
\]

以典型成分 $x = 0.004$（各 0.4 at.\% Ti 和 C）为例：
\[
  \Delta T_R \approx (-25) \times 0.4 + (-60) \times 0.4 \approx -34\ \text{K}
\]
即液相线下移约 34 K。

\subsection{初始浓度}

\begin{table}[h!]
\centering
\caption{初始浓度参数}
\begin{tabular}{llll}
\toprule
符号 & 含义 & 典型值 & 单位 \\
\midrule
$C_1^0$ & Ti 初始浓度 & 0.4 at.\%（可调范围 0.2–1.0） & at.\% \\
$C_2^0$ & C 初始浓度 & 0.4 at.\%（通常与 Ti 等摩尔） & at.\% \\
\bottomrule
\end{tabular}
\end{table}

\subsection{扩散动力学参数}

\begin{table}[h!]
\centering
\caption{扩散动力学参数}
\begin{tabular}{llllll}
\toprule
符号 & 参数含义 & Ti & C & 单位 & 数据来源 \\
\midrule
$D_L$ & 液相扩散系数 & $3\times10^{-9}$ & $8\times10^{-9}$ & m$^2$/s & [13][16] \\
$D_S$ & 固相扩散系数 & $\approx 0$ (Scheil) & $10^{-14} \sim 10^{-13}$ & m$^2$/s & DICTRA (MOBHEA2) [9] \\
\bottomrule
\end{tabular}
\end{table}

\noindent\textbf{注：} 由于 CoCrNi 迟缓扩散效应，固态反向扩散项 $D_S$ 对于 Ti 可安全忽略（Scheil-Gulliver 条件）。

\subsection{热物理参数}

\begin{table}[h!]
\centering
\caption{其他热物理参数}
\begin{tabular}{lllll}
\toprule
符号 & 参数 & 推荐值 & 单位 & 数据来源 \\
\midrule
$L$ & 熔化潜热 & $2.35 \times 10^5$ & J/kg & 实验测定 [10] \\
$\rho$ & 密度 & $8100 \sim 8300$ & kg/m$^3$ & [21] \\
$c_p$ & 比热容 & $530 \sim 560$ & J/(kg$\cdot$K) & [10] \\
$\sigma_{SL}$ & 固液界面能 & $0.18 \sim 0.25$ & J/m$^2$ & 相场模拟反演 [17] \\
$\Gamma$ & Gibbs-Thomson 系数 & $(1\sim3)\times10^{-7}$ & K$\cdot$m & [18] \\
$\varepsilon_4$ & 界面能各向异性强度 & $0.03 \sim 0.05$ & — & [13] \\
\bottomrule
\end{tabular}
\end{table}

\subsection{形核参数（针对 CET 建模）}

\begin{table}[h!]
\centering
\caption{形核参数}
\begin{tabular}{llll}
\toprule
参数 & 含义 & 推荐值 & 单位 \\
\midrule
$n_{\max,v}$ & 体积形核密度（活性基底总数） & $10^{11} \sim 10^{13}$ & m$^{-3}$ \\
$\Delta T_{N,v}$ & 体积形核激活阈值（平均） & $1.5 \sim 4$ & K \\
$\sigma_{\Delta T,v}$ & 形核分布宽度（标准差） & $0.5 \sim 1$ & K \\
$n_{\max,s}$ & 表面形核密度（壁面活性基底） & $10^{12} \sim 10^{14}$ & m$^{-2}$ \\
$\Delta T_{N,s}$ & 表面形核过冷度 & $\approx 0.5$ & K \\
\bottomrule
\end{tabular}
\end{table}

\subsection{参数使用说明}

将以上参数代入文档主体公式时，对应关系如下：
\begin{itemize}
  \item $m_1 = -25\ \text{K/at.\%}$（Ti 液相线斜率，建议值）
  \item $m_2 = -60\ \text{K/at.\%}$（C 液相线斜率，建议值）
  \item $k_1 = 0.08$（Ti 平衡分配系数，建议值）
  \item $k_2 = 0.05$（C 平衡分配系数，建议值）
  \item $T_0 = 1610\ \text{K}$（CoCrNi 基体液相线温度）
  \item $\Delta T_R$：用实际液相线修正（由 CALPHAD 给出，或用上述近似计算）
\end{itemize}

求解结果 $C_1^{eq}$、$C_2^{eq}$ 代表固液界面处 Ti 和 C 的平衡液相浓度；$\Delta f_S$ 代表每步固相分数增量，反映凝固前沿推进速率。

\bigskip
\noindent\textit{参数来源：PDF 报告《3D CA Simulation of Dendrite Growth in CoCrNi-Ti-C Alloys》，2026年3月}

\end{document}
